\chapter{Discussion and Conclusion}
\section*{Discussion}
\addcontentsline{toc}{section}{Discussions}

This thesis presented two complementary works originally motivated by challenges in evaluating GRN reconstruction algorithms, specifically from single-cell transcriptomics data. We first presented GRouNdGAN, a causal implicit generative model inspired by CausalGAN \cite{kocaoglu2017causalgan}, a simulator from the computer vision domain. This method accepts an scRNA-seq reference and a set of user-supplied regulatory TF-gene interactions and adversarially trains a generator neural network consistent with both. A primary strength of our approach is that these interactions do not need to reflect the true causal relationships of the reference (as those are unknown). But, empirical and theoretical guarantees demonstrate that they will collectively form the underlying causal graph of the synthetic data, which can serve as ground truth when inferring GRNs. The model’s ability to disrupt interactions present in the reference data but not imposed during simulation is particularly essential, as it avoids false-positive inference of such edges. We used this usecase to benchmark various GRN inference methods, obtaining results that aligned with experimental and curated benchmarks. 


Our method overcomes major limitations faced by GRN-guided simulators, such as GNW \cite{schaffter2011genenetweaver}, BoolODE \cite{pratapa2020benchmarking}, SERGIO \cite{dibaeinia2020sergio}, and scMultiSim \cite{li2025scmultisim}, which follow kinetic modeling. As de novo methods, they inherently do not mimic experimental data. As such, they rely on post-hoc matching of dataset-, gene-, or cell-level statistical measures, which do not capture the complexities of transcriptomics data, sometimes even passing on the burden to the user. This process even risks disrupting already-imposed interactions. Moreover, being rule-based, they must make coarse assumptions regarding gene expression and TF coregulation. Being a deep learning model, GRouNdGAN makes no such simplifying assumptions and preserves all properties of its reference dataset, including gene identities, cell type markers, lineage trajectories, and pseudo-time orderings, automatically.

As opposed to existing simulators that can only sample from observational data available in their training dataset, GRouNdGAN’s can sample from interventional distributions to predict the effect of TF perturbations (knockout, overexpression, etc.) in silico. Since perturbations are carried out on the same set of cells, keeping all other simulation parameters constant, it is possible to form matched case/control data for validating counterfactual inference methods. Creating such data is not possible on a large scale experimentally, as cells are lysed to extract RNA content when sequenced. 

Beyond methodological innovation, this thesis emphasizes the practical utility of our method. We have published a database of simulated datasets with various GRNs and reference datasets to be used by researchers without the need for model training. These datasets have since supported the development and evaluation of new GRN inference methods \cite{zandigohar2025scregulate, afonja2024llm4grn}. To make our method, which is comprised of multi-step preprocessing, training, and evaluation pipelines, easily reproducible and simple to run and use for all purposes discussed, we developed a toolkit. This toolkit also includes metrics for assessing the quality of simulations and inferred GRNs. Alongside, we published a paper that provides procedural instructions on using the toolkit, as well as additional guidance on equipment setup, installation, input preparation, hyperparameters, training, monitoring, timing, and troubleshooting common issues. This work also included studies on simulation and benchmarking using cell-type-specific GRNs and robustness and stability analyses through data augmentation. We expect these efforts to streamline GRN method development.

Several findings from our work warrant further discussion. As it currently stands, our method has a few constraints that can be relaxed in the future. First, the incorporation of the LSN layer was shown to boost realism in our ablation studies. However, its effect is less noticeable when simulating fewer genes (1000 instead of 20000, as was the case with scGAN \cite{marouf2020realistic}). This may suggest that certain architectural components are more beneficial under specific data regimes, which would require their own ablation. Second, while GRouNdGAN currently enforces single-layer GRNs (bipartite TF-to-gene graphs), it is trivially possible to extend the method to enforce GRNs with depth > 1 on the generator. Cycles could be incorporated by unfolding the graph and following a multi-step training approach to generate genes layer by layer. MOSim \cite{monzo2025mosim}, another network-guided simulator, allows multi-layer GRNs. 

Another important consideration is that we observed that simulators that do not explicitly incorporate GRNs lose the regulatory information of the reference, despite scoring well on realism metrics. GRouNdGAN can also be used for general-purpose simulation tasks. Despite being semi-mechanistic and causally constrained, it outperforms state-of-the-art simulators in realistic generation when provided with a GRN consistent with the reference, as it uses it as additional domain knowledge.

This work also underscores broader implications for computational biology. GRouNdGAN can also be extended to condition on additional labels such as cell type, cluster, or location. Our network-guided simulation framework can be adapted to generate data from other modalities (e.g., ATAC-seq and ChIP-Seq) and networks (e.g., protein-protein interaction and signaling networks) and to inspire the development of causal generative models for biology, based on newer models, for example, diffusion and transformers. Currently, GRouNdGAN requires an input GRN, either supplied by the user or inferred using tools like GRNBoost2 \cite{moerman2019grnboost2}. A promising direction would be to integrate GRN learning into the generative process itself, potentially by enforcing sparsity constraints to encourage the discovery of causal structures, similar to DAG-GAN \cite{gao2021dag}. 

\section*{Conclusion}
\addcontentsline{toc}{section}{Conclusion}
In this thesis, we developed GRouNdGAN and made it accessible to researchers through a toolkit, and formalized the framework through a protocol publication. Our method's objectives are two-fold: simulating realistic single-cell data while causally imposing user-provided regulatory interactions. Our key contributions include: 

\begin{enumerate}
\item Demonstrating that GRouNdGAN-generated data can serve as ground truth, effectively bridging the gap between synthetic and experimental benchmarks of GRN inference methods.
\item Providing a platform for in silico matched case/control TF perturbations.
\item Developing publicly available datasets of GRN ground truths and a toolkit to ensure reproducibility, accelerate GRN inference method development, and facilitate adoption.
\end{enumerate}



% talk about ill-posed Observational vs. Interventional
% 3. Biological Latency and "Snapshots"
%Transcription is a dynamic, time-dependent process.

%The Delay: A Transcription Factor (TF) protein is produced, enters the nucleus, binds to DNA, and then the target gene is transcribed.
%
%The Mismatch: In scRNA-seq, we measure mRNA levels, not protein activity. Since there is a significant time lag and complex translation regulation between mRNA and functional protein, the "cause" (TF activity) and the "effect" (target gene expression) often aren't captured in the same state.